\hypertarget{conventions-and-finite-packet-space}{%
\section{Conventions and finite packet
space}\label{conventions-and-finite-packet-space}}

We use natural units \(\hbar=c=1\). Space is the flat unit torus

\[
\mathbb T^3=(\mathbb R/\mathbb Z)^3,
\]

equipped with normalized Haar measure \(dx\), so
\(\int_{\mathbb T^3}1\,dx=1\). Integration of a periodic representative
may equivalently be taken over \([0,1)^3\) with Lebesgue measure. A
Fourier mode is \(n=(n_1,n_2,n_3)\in\mathbb Z^3\), its momentum is
\(p(n)=2\pi n\), and its scalar plane wave is

\[
e_n(x)=e^{2\pi i n\cdot x}=e^{i p(n)\cdot x}.
\]

Components of a spinor are numbered \(0,1,2,3\), in the exact order of
the source representation. For \(v,w\in\mathbb C^4\),

\[
\langle v,w\rangle=\sum_{r=0}^{3}\overline{v_r}w_r,
\qquad
\|v\|_2^2=\sum_{r=0}^{3}|v_r|^2.
\]

The inner product is conjugate-linear in its first argument and linear
in its second. The adjoint is therefore the conjugate transpose. Unless
indicated otherwise, a spinor norm is this Euclidean norm and an
operator norm is its induced operator norm. The field pairing is

\[
\langle f,g\rangle_{L^2}
=\int_{\mathbb T^3}\langle f(x),g(x)\rangle\,dx
\]

on \(L^2(\mathbb T^3;\mathbb C^4)\). Fix a finite set
\(S\subset\mathbb Z^3\). A packet with amplitudes \(a_n\in\mathbb C^4\)
is

\[
f_a(x)=\sum_{n\in S}e_n(x)a_n.
\]

The associated coefficient space is the finite Hilbert direct sum

\[
\mathcal K_S=\bigoplus_{n\in S}\mathbb C^4,
\qquad
\|a\|_{\mathcal K_S}^2=\sum_{n\in S}\|a_n\|_2^2.
\]

The packet space \(\mathcal F_S=\{f_a:a\in\mathcal K_S\}\) will be
identified isometrically with \(\mathcal K_S\). Its elements are smooth
periodic functions, so the differential and centered-translation
calculations below raise no regularity question. Empty support is
allowed: both spaces then contain only zero and all packet identities
reduce to \(0=0\).

The parameters \(m,h,t\) are arbitrary real numbers. The mathematics
does not require \(m\ge0\); the usual physical rest-mass interpretation
uses that additional convention. The shell value \(E\) is always a
nonnegative square root, even when \(m<0\). Mesh limits in the checked
statements are \(h\to0^+\). Some explicit estimates proved here hold for
every real \(h\), since the centered symbol is even in \(h\).

\hypertarget{the-coordinate-clifford-representation}{%
\section{The coordinate Clifford
representation}\label{the-coordinate-clifford-representation}}

The spatial generators and mass generator are

\[
\alpha_x=
\begin{pmatrix}
0&0&0&1\\
0&0&1&0\\
0&1&0&0\\
1&0&0&0
\end{pmatrix},
\qquad
\alpha_y=
\begin{pmatrix}
0&0&0&-i\\
0&0&i&0\\
0&-i&0&0\\
i&0&0&0
\end{pmatrix},
\]

\[
\alpha_z=
\begin{pmatrix}
0&0&1&0\\
0&0&0&-1\\
1&0&0&0\\
0&-1&0&0
\end{pmatrix},
\qquad
\beta=
\begin{pmatrix}
1&0&0&0\\
0&1&0&0\\
0&0&-1&0\\
0&0&0&-1
\end{pmatrix}.
\]

In particular, \(\alpha_yv=(-iv_3,iv_2,-iv_1,iv_0)\); this fixes a sign
that cannot be changed silently in a worked example. The matrices act on
column spinors. Grouping the first two and last two components gives
their block form without any coordinate permutation:

\[
\alpha_j=
\begin{pmatrix}0&\sigma_j\\\sigma_j&0\end{pmatrix},
\qquad
\beta=\begin{pmatrix}I_2&0\\0&-I_2\end{pmatrix},
\]

where

\[
\sigma_x=\begin{pmatrix}0&1\\1&0\end{pmatrix},
\quad
\sigma_y=\begin{pmatrix}0&-i\\i&0\end{pmatrix},
\quad
\sigma_z=\begin{pmatrix}1&0\\0&-1\end{pmatrix}.
\]

\needspace{4\baselineskip}
\textbf{Proposition PH01 (Hermitian Clifford relations).} The four
displayed matrices are Hermitian involutions. Distinct spatial
generators anticommute, and each spatial generator anticommutes with
\(\beta\):

\[
\alpha_j^*=\alpha_j,\quad \beta^*=\beta,\quad
\alpha_j^2=\beta^2=I_4,
\]

\[
\alpha_j\alpha_k+\alpha_k\alpha_j=0\quad(j\ne k),
\qquad
\alpha_j\beta+\beta\alpha_j=0.
\]

\emph{Proof.} The real diagonal entries of \(\beta\) are
\(1,1,-1,-1\), so its adjoint and square have the stated forms. Each
\(\sigma_j\) is Hermitian and squares to \(I_2\), by direct
multiplication. Their products are

\[
\sigma_x\sigma_y=i\sigma_z,\qquad
\sigma_y\sigma_z=i\sigma_x,\qquad
\sigma_z\sigma_x=i\sigma_y,
\]

with reversed products equal to their negatives. For example,

\[
\begin{pmatrix}0&1\\1&0\end{pmatrix}
\begin{pmatrix}0&-i\\i&0\end{pmatrix}
=\begin{pmatrix}i&0\\0&-i\end{pmatrix}=i\sigma_z.
\]

The other two follow by the same two-by-two multiplication; reversal
follows by taking adjoints. The block form gives

\[
\alpha_j^*=
\begin{pmatrix}0&\sigma_j^*\\\sigma_j^*&0\end{pmatrix}
=\alpha_j,\qquad
\alpha_j\alpha_k=
\begin{pmatrix}\sigma_j\sigma_k&0\\0&\sigma_j\sigma_k\end{pmatrix}.
\]

Thus \(\alpha_j^2=I_4\) and the spatial anticommutators vanish. Finally,

\[
\alpha_j\beta=
\begin{pmatrix}0&-\sigma_j\\\sigma_j&0\end{pmatrix},
\qquad
\beta\alpha_j=
\begin{pmatrix}0&\sigma_j\\-\sigma_j&0\end{pmatrix},
\]

whose sum is zero. Every identity is an identity of the exact source
matrices. \(\square\)

\hypertarget{continuum-and-centered-mode-hamiltonians}{%
\section{Continuum and centered mode
Hamiltonians}\label{continuum-and-centered-mode-hamiltonians}}

For \(p\in\mathbb R^3\), define

\[
H(m,p)=m\beta+p_x\alpha_x+p_y\alpha_y+p_z\alpha_z.
\]

Its coordinates are

\[
H(m,p)=
\begin{pmatrix}
m&0&p_z&p_x-ip_y\\
0&m&p_x+ip_y&-p_z\\
p_z&p_x-ip_y&-m&0\\
p_x+ip_y&-p_z&0&-m
\end{pmatrix}.
\]

On mode \(n\), the continuum Hamiltonian is \(H_n=H(m,p(n))\). It is the
symbol of the differential expression

\[
-i\sum_{j=x,y,z}\alpha_j\partial_j+m\beta
\]

on the finite packet space: \(\partial_je_n=ip_j(n)e_n\), so applying
the expression to \(e_na_n\) gives \(e_nH_na_n\). For real spacing
\(h\), set

\[
q_j(h,p)=
\begin{cases}
\sin(hp_j)/h,&h\ne0,\\
p_j,&h=0.
\end{cases}
\]

Write \(q_h(n)=q(h,p(n))\) and \(H_{h,n}=H(m,q_h(n))\). The function
\(q_j\) is even in \(h\), and its value at zero is justified by
\(\sin z/z\to1\). For \(h\ne0\), a centered translation difference on a
periodic smooth function is

\[
(D_{h,j}f)(x)=\frac{f(x+he_j)-f(x-he_j)}{2h}.
\]

Here \(e_j\) is the j-th spatial unit vector, distinct from the Fourier
characters \(e_n\) and spinor vectors. Its action on a plane wave is
exactly

\[
D_{h,j}e_n
=e_n\frac{e^{ihp_j(n)}-e^{-ihp_j(n)}}{2h}
=i\frac{\sin(hp_j(n))}{h}e_n.
\]

The centered differential expression \(-i\sum_j\alpha_jD_{h,j}+m\beta\)
therefore has symbol \(H_{h,n}\). Its signs follow from \(-i\cdot i=1\).
At \(h=0\), the symbol equals \(H_n\) by definition. These formulas take
place on the continuum torus and use no sampled spatial data.

\needspace{4\baselineskip}
\textbf{Proposition PH02 (Hermiticity and the mass shell).} For every
real \(m\) and \(p\in\mathbb R^3\),

\[
H(m,p)^*=H(m,p),\qquad
H(m,p)^2=(m^2+|p|^2)I_4.
\]

In particular, \(H_n,H_{h,n}\) are Hermitian for all real \(m,h\), and

\[
H_n^2=E_n^2I_4,\quad E_n=\sqrt{m^2+|p(n)|^2},
\]

\[
H_{h,n}^2=E_{h,n}^2I_4,\quad
E_{h,n}=\sqrt{m^2+|q_h(n)|^2}.
\]

\emph{Proof.} Adjoint is additive, and a real multiple of a Hermitian
matrix is Hermitian. PH01 therefore proves the first identity. Expanding
the square gives the square terms \(m^2\beta^2+\sum_jp_j^2\alpha_j^2\),
the terms \(mp_j(\beta\alpha_j+\alpha_j\beta)\), and the terms
\(p_jp_k(\alpha_j\alpha_k+\alpha_k\alpha_j)\) for \(j<k\). The
anticommutators vanish and all four squares are \(I_4\), giving
\((m^2+\sum_jp_j^2)I_4\). Substitution proves both mode statements.
\(\square\)

Hermiticity here is a statement about matrices acting on all of
\(\mathbb C^4\). Every linear map on this finite space is bounded and
everywhere defined. This is not a construction of the adjoint or domain
of an unbounded differential operator on the full field space.

\hypertarget{spectral-projectors-and-exact-mode-evolution}{%
\section{Spectral projectors and exact mode
evolution}\label{spectral-projectors-and-exact-mode-evolution}}

Let \(H=H(m,p)\) and \(E=\sqrt{m^2+|p|^2}\). Define the exact propagator

\[
U_H(t)=e^{-itH}=\sum_{k=0}^{\infty}\frac{(-itH)^k}{k!}.
\]

The series converges absolutely in operator norm, being dominated by
\(\sum_k(|t|\|H\|)^k/k!\). This is exact continuous-time evolution, with
no time step or series truncation.

\needspace{4\baselineskip}
\textbf{Proposition PH03 (spectral splitting and closed propagator).} If
\(E>0\), define

\[
P_+=\frac12\left(I_4+\frac HE\right),\qquad
P_-=\frac12\left(I_4-\frac HE\right).
\]

Then \(P_\pm^*=P_\pm\), \(P_\pm^2=P_\pm\), \(P_+P_-=P_-P_+=0\),
\(P_++P_-=I_4\), and \(HP_\pm=\pm EP_\pm\). Both projectors have rank
two, and

\[
U_H(t)=e^{-itE}P_++e^{itE}P_-
=\cos(tE)I_4-i\frac{\sin(tE)}{E}H.
\]

If \(E=0\), then \(H=0\) and \(U_H(t)=I_4\). Equivalently, the closed
expression is valid with the continuous convention \(\sin(tE)/E=t\) at
\(E=0\).

\emph{Proof.} For \(E>0\), the Hermitian matrix \(K=H/E\) satisfies
\(K^2=I_4\). Thus

\[
P_\pm^2=\tfrac14(I_4\pm2K+K^2)=\tfrac12(I_4\pm K)=P_\pm,
\]

and \(P_+P_-=\tfrac14(I_4-K^2)=0\); the reversed product is the same.
Their adjoints equal themselves and their sum is \(I_4\). Also
\(K(I_4\pm K)=K\pm I_4=\pm(I_4\pm K)\), proving \(HP_\pm=\pm EP_\pm\).
Each displayed Clifford matrix has trace zero. Hence
\(\operatorname{tr}H=0\) and \(\operatorname{tr}P_\pm=2\). An idempotent
splits the space into its range and kernel, acting as identity and zero
respectively. Its trace is therefore its rank, proving rank two.

The identity \(I_4=P_++P_-\) and the two eigenvalue identities give, by
induction, \(H^k=E^kP_++(-E)^kP_-\) for every \(k\ge0\). Substitute into
the absolutely convergent series to obtain the spectral expression.
Expanding its projectors gives the cosine and sine expression.

When \(E=0\), the sum \(m^2+|p|^2\) of nonnegative squares is zero.
Therefore \(m=0,p=0\), and \(H=0\). Its exponential is \(I_4\). The
quotient \(\sin(tE)/E\) tends to \(t\) at zero; multiplying its filled
value by \(H=0\) gives the same identity. No division by zero is used.
\(\square\)

The proof applies unchanged to \(H_{h,n}\), with energy \(E_{h,n}\). A
positive-energy continuum mode can have \(E_{h,n}=0\) at a nonzero
spacing when \(m=0\) and all centered components vanish. That mode uses
the zero-energy branch. Positivity of the continuum energy does not
license division by the centered energy at every spacing.

\needspace{4\baselineskip}
\textbf{Proposition PH04 (group law, unitarity, and stability).} For
real \(s,t\),

\[
U_H(0)=I_4,\qquad U_H(s+t)=U_H(s)U_H(t),
\]

\[
U_H(-t)=U_H(t)^{-1}=U_H(t)^*,
\qquad U_H(t)^*U_H(t)=U_H(t)U_H(t)^*=I_4.
\]

Consequently,

\[
\langle U_H(t)v,U_H(t)w\rangle=\langle v,w\rangle,\qquad
\|U_H(t)v\|_2=\|v\|_2
\]

for all spinors. These statements hold for both mode Hamiltonians at
every real \(m,h,t\).

\emph{Proof.} When \(E>0\), the spectral form and orthogonality of
\(P_\pm\) give

\[
U_H(s)U_H(t)
=e^{-i(s+t)E}P_++e^{i(s+t)E}P_-=U_H(s+t).
\]

At zero time the sum is \(I_4\). Hermiticity of the projectors and
reality of \(tE\) give \(U_H(t)^*=e^{itE}P_++e^{-itE}P_-=U_H(-t)\).
Setting \(s=-t\) proves the inverse identities. When \(E=0\), every
assertion follows from \(U_H(t)=I_4\).

For any matrix \(A\), expansion and interchange of two finite sums gives
\(\langle Av,w\rangle=\langle v,A^*w\rangle\). Thus

\[
\langle U_H(t)v,U_H(t)w\rangle
=\langle v,U_H(t)^*U_H(t)w\rangle=\langle v,w\rangle.
\]

Take \(w=v\) and then nonnegative square roots to obtain norm
preservation. PH02 supplies the mode cases. \(\square\)

An equivalent bounded-operator argument is represented in the formal source:
\(-itH\) is skew-adjoint, and its exponential is unitary. Our spectral
proof supplies additional explicit projectors and formulas.

\hypertarget{fourier-orthogonality-and-packet-conservation}{%
\section{Fourier orthogonality and packet
conservation}\label{fourier-orthogonality-and-packet-conservation}}

Let \(J_n:\mathbb C^4\to L^2(\mathbb T^3;\mathbb C^4)\) be
\(J_nv=e_nv\).

\needspace{4\baselineskip}
\textbf{Proposition PH05 (mode isometry and finite Parseval identity).}
For modes \(n,k\) and spinors \(v,w\),

\[
\langle J_nv,J_kw\rangle_{L^2}
=\begin{cases}\langle v,w\rangle,&n=k,\\0,&n\ne k.\end{cases}
\]

In particular, \(J_n\) is a linear isometry and

\[
\left\|\sum_{n\in S}J_na_n\right\|_{L^2}^2
=\sum_{n\in S}\|a_n\|_2^2.
\]

\emph{Proof.} The constant spinors factor out of the scalar integral:

\[
\langle J_nv,J_kw\rangle_{L^2}
=\langle v,w\rangle\int_{\mathbb T^3}e^{2\pi i(k-n)\cdot x}\,dx.
\]

The integral is the product of three one-dimensional integrals. For
integer \(r\), the integral of \(e^{2\pi irx}\) over \([0,1)\) is \(1\)
when \(r=0\); for \(r\ne0\), it is \((e^{2\pi ir}-1)/(2\pi ir)=0\). The
product is therefore \(1\) exactly when \(n=k\). Linearity is immediate,
and setting \(k=n,w=v\) proves isometry. Expansion of the squared norm
of the finite sum gives a double sum of pairings. The off-diagonal
pairings vanish and the diagonal pairings are \(\|a_n\|_2^2\).
\(\square\)

Define the evolved packets

\[
f(t,x)=\sum_{n\in S}e_n(x)U_{H_n}(t)a_n,\qquad
f_h(t,x)=\sum_{n\in S}e_n(x)U_{H_{h,n}}(t)a_n.
\]

Both remain in \(\mathcal F_S\). Their operators \(U_S(t),U_{h,S}(t)\)
have the finite block diagonal coefficient actions with blocks
\(U_{H_n}(t),U_{H_{h,n}}(t)\).

\needspace{4\baselineskip}
\textbf{Theorem PH06 (exact finite-packet norm conservation).} For every
real \(m,h,t\), finite \(S\), and coefficient family \(a\),

\[
\|f(t)\|_{L^2}^2=\|f_h(t)\|_{L^2}^2
=\sum_{n\in S}\|a_n\|_2^2=\|f_a\|_{L^2}^2.
\]

Both packet evolutions preserve the \(L^2\) inner product between two
packets in \(\mathcal F_S\), and each forms a one-parameter unitary
group there.

\emph{Proof.} Apply PH05 to the evolved coefficients and PH04 on each
mode:

\[
\|f_h(t)\|_{L^2}^2
=\sum_{n\in S}\|U_{H_{h,n}}(t)a_n\|_2^2
=\sum_{n\in S}\|a_n\|_2^2.
\]

The continuum argument is identical. For coefficient families \(a,b\),
the same two propositions give

\[
\langle f_{h,a}(t),f_{h,b}(t)\rangle_{L^2}
=\sum_{n\in S}
\langle U_{H_{h,n}}(t)a_n,U_{H_{h,n}}(t)b_n\rangle
=\sum_{n\in S}\langle a_n,b_n\rangle.
\]

Composition of coefficient blocks proves the packet group law, and
inverse time proves surjectivity. This establishes unitarity on
\(\mathcal F_S\) using only finite sums. \(\square\)

The quantity conserved is squared Hilbert norm, or total probability for
a normalized state. The parameter \(m\), the positive shell value
\(E_n\), and Hamiltonian expectation are different quantities. Calling
squared norm a packet mass does not identify it with the rest mass or
physical Hamiltonian energy.

\hypertarget{generator-consistency-and-dispersion-error}{%
\section{Generator consistency and dispersion
error}\label{generator-consistency-and-dispersion-error}}

The mode Hamiltonian difference contains no mass term:

\[
\Delta_{h,n}:=H_{h,n}-H_n
=\sum_{j=x,y,z}\delta_j(h,n)\alpha_j,
\qquad
\delta_j(h,n)=q_j(h,p(n))-p_j(n).
\]

Define

\[
C(n)=\frac16\left(\sum_{j=x,y,z}|p_j(n)|^6\right)^{1/2}.
\]

This constant depends on the fixed mode and increases cubically with
momentum scale. That dependence is essential in a fixed-cutoff result.

\needspace{4\baselineskip}
\textbf{Proposition PH07 (quadratic generator consistency).} For every
real \(h\) and fixed mode \(n\),

\[
\|\Delta_{h,n}\|=|\delta(h,n)|\le h^2C(n).
\]

For every fixed finite packet,

\[
\left\|\sum_{n\in S}e_n\Delta_{h,n}a_n\right\|_{L^2}^2
=\sum_{n\in S}|\delta(h,n)|^2\|a_n\|_2^2,
\]

and

\[
\left\|\sum_{n\in S}e_n\Delta_{h,n}a_n\right\|_{L^2}
\le h^2\left(\sum_{n\in S}C(n)^2\|a_n\|_2^2\right)^{1/2}.
\]

\emph{Proof.} First prove the scalar estimate without a sign
restriction on \(h\). The fundamental theorem of calculus gives
\(|\sin u|\le|u|\), since \(\sin u=\int_0^u\cos v\,dv\) and
\(|\cos v|\le1\). Therefore

\[
|1-\cos u|=2\sin^2(u/2)\le u^2/2.
\]

For \(z\ge0\), integrate this bound in
\(\sin z-z=\int_0^z(\cos u-1)\,du\) to obtain \(|\sin z-z|\le z^3/6\).
Oddness of the residual proves \(|\sin z-z|\le|z|^3/6\) for negative
\(z\). For \(h\ne0\), substitute \(z=hp_j(n)\) and divide by \(|h|\):

\[
|\delta_j(h,n)|\le \frac{h^2|p_j(n)|^3}{6}.
\]

For \(h=0\), the residual is zero by definition and the bound still
holds. Squaring and adding proves \(|\delta(h,n)|\le h^2C(n)\).

The Hermitian matrix \(\Delta_{h,n}\) satisfies
\(\Delta_{h,n}^2=|\delta(h,n)|^2I_4\) by PH01. Hence

\[
\|\Delta_{h,n}v\|_2^2
=\langle v,\Delta_{h,n}^*\Delta_{h,n}v\rangle
=|\delta(h,n)|^2\|v\|_2^2.
\]

Its induced operator norm is exactly \(|\delta(h,n)|\): the upper bound
holds on every vector and equality holds on every unit vector. Apply
PH05 to the residual coefficients for the packet equality, substitute
the scalar bounds, and take nonnegative square roots. \(\square\)

The checked generator results state corresponding \(O(h^2)\) estimates
as \(h\to0^+\). The explicit constant and all-real inequality above are
traditional deductions. Precisely, \(r(h)=O(h^2)\) means that some
\(M,\varepsilon>0\) satisfy \(\|r(h)\|\le Mh^2\) for
\(0<h<\varepsilon\). Fixed support and amplitudes permit finite
summation of such bounds; no estimate uniform over all integer modes
follows.

The reverse triangle inequality in \(\mathbb R^4\), applied to
\((m,q_h(n))\) and \((m,p(n))\), gives

\[
|E_{h,n}-E_n|\le|q_h(n)-p(n)|\le h^2C(n).
\]

Moreover, \(|\sin(hp_j)/h|\le|p_j|\) for \(h\ne0\), and the filled value
is equal at zero. Thus \(0\le E_{h,n}\le E_n\). For a fixed mode,
Taylor's formula gives

\[
q_j(h,p(n))=p_j(n)-\frac{h^2p_j(n)^3}{6}+O(h^4),
\]

\[
E_{h,n}^2=E_n^2-\frac{h^2}{3}\sum_jp_j(n)^4+O(h^4).
\]

For clarity, Taylor's theorem through degree four has zero even
coefficients for sine, and its fifth-derivative remainder is at most
\(|hp_j|^5/120\). Division by \(|h|\) gives \(O(h^4)\); squaring the
three expansions and summing preserves that order for a fixed mode. If
\(E_n>0\), the square root is smooth near \(E_n^2\), so

\[
E_{h,n}=E_n-\frac{h^2}{6E_n}\sum_jp_j(n)^4+O(h^4).
\]

If \(E_n=0\), then \(m=0,n=0\) and \(E_{h,n}=0\) exactly for every
\(h\). The expansion containing \(1/E_n\) is not used.

For fixed time, the phase error on either energy branch obeys

\[
|e^{\pm itE_{h,n}}-e^{\pm itE_n}|
\le |t|\,|E_{h,n}-E_n|,
\]

by integrating the derivative of \(s\mapsto e^{\pm its}\) between the
two real energies. Phase comparison alone does not compare full spinor
evolution: the eigenspaces also depend on the kinetic vector. The
operator argument below controls both effects.

\needspace{10\baselineskip}
\atlasfigure{04_dispersion}{The continuum and centered shell branches along one momentum axis, with $m=1$, $h=1/12$. Marked points are integer torus modes; intermediate points use the real-symbol extension.}

\hypertarget{fixed-time-propagator-consistency}{%
\section{Fixed-time propagator
consistency}\label{fixed-time-propagator-consistency}}

\needspace{4\baselineskip}
\textbf{Theorem PH08 (fixed-mode and fixed-packet convergence).} For
fixed real \(m,t\) and fixed \(n\),

\[
U_{H_{h,n}}(t)-U_{H_n}(t)=O(h^2)\quad(h\to0^+).
\]

For fixed finite \(S\) and amplitudes \(a\),

\[
f_h(t)-f(t)=O(h^2)
\quad\hbox{in }L^2(\mathbb T^3;\mathbb C^4),
\]

and in particular \(\|f_h(t)-f(t)\|_{L^2}\to0\). The exact residual
identity is

\[
\|f_h(t)-f(t)\|_{L^2}^2
=\sum_{n\in S}\|(U_{H_{h,n}}(t)-U_{H_n}(t))a_n\|_2^2.
\]

\emph{Proof.} Give a fixed-time proof independent of the later
compact-time extension. For arbitrary matrices \(X,Y\), telescoping
products give

\[
X^k-Y^k=\sum_{r=0}^{k-1}X^{k-1-r}(X-Y)Y^r
\quad(k\ge1).
\]

No commutation is assumed. Expanding the summands gives successive
differences \(X^{k-r}Y^r-X^{k-1-r}Y^{r+1}\), which cancel. If
\(\|X\|,\|Y\|\le R\), then

\[
\|X^k-Y^k\|\le kR^{k-1}\|X-Y\|.
\]

Subtract the absolutely convergent exponential series and sum:

\[
\|e^X-e^Y\|
\le\|X-Y\|\sum_{k=1}^{\infty}\frac{kR^{k-1}}{k!}
=e^R\|X-Y\|.
\]

PH02 gives \(\|H_n\|=E_n,\|H_{h,n}\|=E_{h,n}\), by the norm-squared
argument of PH07. Since \(E_{h,n}\le E_n\), take \(X=-itH_{h,n}\),
\(Y=-itH_n\), and \(R=|t|E_n\). The result is

\[
\|U_{H_{h,n}}(t)-U_{H_n}(t)\|
\le |t|e^{|t|E_n}h^2C(n).
\]

This proves the modewise order, including \(t=0\). The difference of
packets is the finite sum of the mode residuals, so PH05 gives its exact
residual identity. Thus

\[
\|f_h(t)-f(t)\|_{L^2}
\le h^2\left(
\sum_{n\in S}|t|^2e^{2|t|E_n}C(n)^2\|a_n\|_2^2
\right)^{1/2}.
\]

The constant is finite for fixed \(t,S,a\), proving the order and
convergence. \(\square\)

The formal algebraic-coordinate proof obtains this order from continuity and local
differentiability of the Banach-algebra exponential. Its operator norm
belongs to the algebraic coordinate spinor type. The coefficient
conversion in its packet definitions is continuous, so its checked
torus-\(L^2\) convergence is genuine. Our Euclidean PH08 also has the
direct proof above. The following bridge explains why its coordinates
agree with the separately defined Hilbert propagators.

\hypertarget{the-algebraic-to-hilbert-coordinate-bridge}{%
\section{The algebraic-to-Hilbert coordinate
bridge}\label{the-algebraic-to-hilbert-coordinate-bridge}}

Let \(\mathcal A\) be the coordinate space \(\{0,1,2,3\}\to\mathbb C\)
with maximum norm \(\|v\|_\infty=\max_r|v_r|\), and let \(\mathcal H\)
be the same coordinates with the Euclidean inner product. The map

\[
T:\mathcal A\longrightarrow\mathcal H,\qquad (Tv)_r=v_r
\]

is complex-linear and bijective; its inverse also retains the
coordinates. It is not generally an isometry between these norms. The
inequalities

\[
\|v\|_\infty\le\|Tv\|_2\le2\|v\|_\infty
\]

show that both directions are continuous. The first follows because each
coordinate square is bounded by the sum; the second because each of four
squares is at most \(\|v\|_\infty^2\).

\needspace{4\baselineskip}
\textbf{Proposition PH09 (coordinate conjugacy of generators and
exponentials).} Let \(A\) be a complex-linear endomorphism of
\(\mathcal A\), and let \(A^{\mathcal H}=TAT^{-1}\). For every complex
\(z\),

\[
e^{zA^{\mathcal H}}=Te^{zA}T^{-1}.
\]

For the source Clifford actions, the continuum and centered Hilbert
generators are exactly these conjugates of their algebraic generators.
Their exact propagators therefore have the same coordinates after
applying \(T\). The algebraic-coordinate packet, transported coefficientwise by
\(T\), agrees with the Hilbert packet for the corresponding
amplitudes.

\emph{Proof.} Induction gives \((TAT^{-1})^k=TA^kT^{-1}\). For
\(k=0\), use \(TT^{-1}=I\); for the induction step, cancel the adjacent
\(T^{-1}T\). Every exponential partial sum therefore satisfies

\[
\sum_{k=0}^{N}\frac{z^k(A^{\mathcal H})^k}{k!}
=T\left(\sum_{k=0}^{N}\frac{z^kA^k}{k!}\right)T^{-1}.
\]

The exponential series converges in both operator norms. Conjugation is
continuous: writing induced norms in the relevant spaces gives

\[
\|TBT^{-1}\|_2\le\|T\|\,\|B\|_\infty\,\|T^{-1}\|
\le2\|B\|_\infty.
\]

Pass to the limit to obtain the identity.

The source's transported Clifford action is defined by
\(v\mapsto T(A(T^{-1}v))\). The same is true for the displayed diagonal
beta action. Conjugation preserves real coefficients and finite sums,
proving the Hamiltonian equality. Apply the exponential identity with
\(z=-it\). Substituting each coefficient equality into the finite
Fourier sum then proves equality of the packets. \(\square\)

The norm bound also transports quadratic operator errors between
realizations. The bridge is a traditional proof in this paper. The
checked source contains inverse coordinate identities and the pairing
identification, but no complete generator-and-exponential conjugacy
theorem. PH09's conjugacy itself remains an ordinary proof in this text.
Blocks 44 and 45 separately establish Hilbert stability and consistency
directly.

\hypertarget{an-exact-three-mode-packet}{%
\section{An exact three-mode packet}\label{an-exact-three-mode-packet}}

\needspace{4\baselineskip}
\textbf{Example PH10 (axis and diagonal modes at \(h=1/12\)).} Set
\(m=1\), \(h=1/12\), and

\[
S=\{n_0,n_A,n_D\},\qquad
n_0=(0,0,0),\quad n_A=(1,0,0),\quad n_D=(1,1,1).
\]

Let \(e_0=(1,0,0,0)^T\) denote the first spinor basis vector, and choose

\[
a_{n_0}=\frac1{\sqrt2}e_0,\qquad
a_{n_A}=\frac12e_0,\qquad
a_{n_D}=\frac12e_0.
\]

Their squared norms are \(1/2,1/4,1/4\), so the packet is normalized.
Its continuum and centered momenta are

\needspace{8\baselineskip}
\begin{center}
\begin{tabular}{@{}lll@{}}
\toprule
Mode & \(p(n)\) & \(q_h(n)\) \\
\midrule

\(n_0\) & \((0,0,0)\) & \((0,0,0)\) \\
\(n_A\) & \((2\pi,0,0)\) & \((6,0,0)\) \\
\(n_D\) & \((2\pi,2\pi,2\pi)\) & \((6,6,6)\) \\
\bottomrule
\end{tabular}
\end{center}

Each nonzero component has \(hp_j=\pi/6\), so
\(\sin(hp_j)/h=(1/2)/(1/12)=6\). Thus

\needspace{8\baselineskip}
\begin{center}
\begin{tabular}{@{}lll@{}}
\toprule
Mode & \(E_n\) & \(E_{h,n}\) \\
\midrule

\(n_0\) & \(1\) & \(1\) \\
\(n_A\) & \(\sqrt{1+4\pi^2}\) & \(\sqrt{37}\) \\
\(n_D\) & \(\sqrt{1+12\pi^2}\) & \(\sqrt{109}\) \\
\bottomrule
\end{tabular}
\end{center}

These are exact symbolic equalities. The centered component \(6\) is
distinct from the integer mode label \(1\) and continuum momentum
\(2\pi\).

For any real kinetic vector \(r\), the exact matrix gives

\[
H(1,r)e_0=(1,0,r_z,r_x+ir_y)^T.
\]

Set \(E_r=\sqrt{1+|r|^2}\), \(c_r(t)=\cos(E_rt)\), and
\(s_r(t)=\sin(E_rt)/E_r\). PH03 gives

\[
u_r(t)=U_{H(1,r)}(t)e_0=
\begin{pmatrix}
c_r(t)-is_r(t)\\
0\\
-ir_zs_r(t)\\
(r_y-ir_x)s_r(t)
\end{pmatrix}.
\]

In particular, \(u_0(t)=e^{-it}e_0\). For the centered axis and diagonal
modes,

\[
u_A(t)=
\begin{pmatrix}
\cos(\sqrt{37}t)-i\sin(\sqrt{37}t)/\sqrt{37}\\
0\\
0\\
-6i\sin(\sqrt{37}t)/\sqrt{37}
\end{pmatrix},
\]

\[
u_D(t)=
\begin{pmatrix}
\cos(\sqrt{109}t)-i\sin(\sqrt{109}t)/\sqrt{109}\\
0\\
-6i\sin(\sqrt{109}t)/\sqrt{109}\\
6(1-i)\sin(\sqrt{109}t)/\sqrt{109}
\end{pmatrix}.
\]

The complete centered packet is

\[
f_h(t,x)=\frac{e^{-it}}{\sqrt2}e_0
+\frac12e^{2\pi ix_1}u_A(t)
+\frac12e^{2\pi i(x_1+x_2+x_3)}u_D(t).
\]

The continuum packet has the same form with \(u_A,u_D\) replaced by
\(u_{(2\pi,0,0)}\) and \(u_{(2\pi,2\pi,2\pi)}\). At \(t=0\), both equal

\[
f_a(x)=\left(\frac1{\sqrt2}+\frac12e^{2\pi ix_1}
+\frac12e^{2\pi i(x_1+x_2+x_3)}\right)e_0.
\]

The spinor formula verifies stability directly:

\[
\|u_r(t)\|_2^2
=c_r(t)^2+s_r(t)^2(1+r_z^2+r_x^2+r_y^2)
=\cos^2(E_rt)+\sin^2(E_rt)=1.
\]

Here \(|c-is|^2=c^2+s^2\) and \(|r_y-ir_x|^2=r_x^2+r_y^2\), since the
coefficients are real. Fourier orthogonality gives

\[
\|f_h(t)\|_{L^2}^2=\tfrac12+\tfrac14+\tfrac14=1,\qquad
\|f(t)\|_{L^2}^2=1.
\]

This is an integrated statement. The initial pointwise density is not
constant: at \(x=0\), it is \((1+1/\sqrt2)^2=3/2+\sqrt2\). Cross terms
can change local density while their torus integrals vanish.
Normalization does not mean that each point has density one.

The error also has an exact expression. For fixed mode, let \(r=p(n)\),
\(q=q_h(n)\), with the corresponding \(c_r,s_r,c_q,s_q\). Coordinate
subtraction gives

\[
\|u_q(t)-u_r(t)\|_2^2
=(c_q-c_r)^2+(s_q-s_r)^2
+\sum_{j=x,y,z}(q_js_q-r_js_r)^2.
\]

For the selected packet, the zero mode contributes no error, so

\[
\begin{aligned}
\|f_h(t)-f(t)\|_{L^2}^2
={}&\frac14\|u_{(6,0,0)}(t)-u_{(2\pi,0,0)}(t)\|_2^2\\
&+\frac14\|u_{(6,6,6)}(t)-u_{(2\pi,2\pi,2\pi)}(t)\|_2^2.
\end{aligned}
\]

A plot or numerical evaluation illustrates this exact expression;
conservation and normalization have already been proved symbolically.

\needspace{10\baselineskip}
\atlasfigure{06_packet_evolution}{The centered three-mode packet at $m=1$, $h=1/12$ changes local density and spinor composition while keeping total squared norm one. The image evaluates the closed formula at finitely many times.}

\hypertarget{a-finite-matrix-compact-time-extension}{%
\section{A finite-matrix compact-time
extension}\label{a-finite-matrix-compact-time-extension}}

The following proof gives an explicit Duhamel bound, including its sharp coefficient. This is an ordinary mathematical deduction with a complete proof. The formal compact-time convergence argument uses exponential continuity and is identified separately in the verification appendix.

\needspace{4\baselineskip}
\textbf{Lemma (finite Hermitian Duhamel estimate).} Let \(A,B\) be
Hermitian matrices on a nonzero finite-dimensional complex Hilbert
space. For every real \(t\),

\[
\|e^{-itA}-e^{-itB}\|\le |t|\,\|A-B\|.
\]

The same inequality holds trivially on the zero space.

\emph{Proof.} First justify differentiation. On \(|s|\le R\), the
series for \(e^{-isB}\) converges uniformly, and its derivative series
is norm-dominated by

\[
\sum_{k=1}^{\infty}\frac{kR^{k-1}\|B\|^k}{k!}
=\|B\|e^{R\|B\|}.
\]

Termwise differentiation therefore gives

\[
\frac{d}{ds}e^{-isB}=-iB e^{-isB}=-ie^{-isB}B.
\]

The same holds for \(A\). Since \(A\) commutes with every power of
itself, it commutes with \(e^{-iuA}\). Fix \(t\) and define
\(F(s)=e^{-i(t-s)A}e^{-isB}\). The product rule gives

\[
F'(s)=i e^{-i(t-s)A}(A-B)e^{-isB}.
\]

The endpoints are \(F(0)=e^{-itA}\) and \(F(t)=e^{-itB}\). Integrating
entrywise in the finite matrix space gives the oriented-integral
identity

\[
e^{-itB}-e^{-itA}
=i\int_0^t e^{-i(t-s)A}(A-B)e^{-isB}\,ds.
\]

For Hermitian matrices the two exponentials in the integrand are
unitary. This does not rely on the special Clifford algebra: the series
gives \((e^{-iuA})^*=e^{iuA}\), and the Cauchy product of the two
absolutely convergent series gives \(e^{iuA}e^{-iuA}=I\). Indeed, the
coefficient of \(A^k\) for \(k>0\) is
\((iu)^k\sum_{r=0}^{k}\binom{k}{r}(-1)^r/k!=0\), and the constant
coefficient is \(I\). A unitary matrix on a nonzero space has norm one.
Submultiplicativity therefore bounds the integrand norm by \(\|A-B\|\).
For \(t\ge0\), integration gives \(t\|A-B\|\); for \(t<0\), reversing
the interval gives \((-t)\|A-B\|\). This proves the estimate for every
real time. \(\square\)

\needspace{4\baselineskip}
\textbf{Theorem PH11 (time-explicit finite-packet bound).} For every
real \(m,h,t\), fixed mode \(n\), and fixed finite packet,

\[
\|U_{H_{h,n}}(t)-U_{H_n}(t)\|
\le |t|\,|\delta(h,n)|\le |t|h^2C(n),
\]

\[
\begin{aligned}
\|f_h(t)-f(t)\|_{L^2}
&\le |t|\left(\sum_{n\in S}|\delta(h,n)|^2\|a_n\|_2^2\right)^{1/2}\\
&\le |t|h^2\left(\sum_{n\in S}C(n)^2\|a_n\|_2^2\right)^{1/2}.
\end{aligned}
\]

Consequently, for every \(T\ge0\),

\[
\sup_{|t|\le T}\|f_h(t)-f(t)\|_{L^2}
\le Th^2\left(\sum_{n\in S}C(n)^2\|a_n\|_2^2\right)^{1/2},
\]

and evolution converges uniformly on every compact time interval for
fixed support and initial amplitudes.

\emph{Proof.} Apply the Duhamel lemma with \(A=H_{h,n}\), \(B=H_n\).
Both are Hermitian by PH02, and PH07 gives the residual norm and bound.
For each amplitude,

\[
\|(U_{H_{h,n}}(t)-U_{H_n}(t))a_n\|_2
\le |t|\,|\delta(h,n)|\,\|a_n\|_2
\le |t|h^2C(n)\|a_n\|_2.
\]

Square and sum over \(S\), apply the exact residual Parseval identity of
PH08, and take the nonnegative square root. The middle coefficient is
exactly the initial generator-residual norm in PH07 because these
Clifford differences satisfy
\(\Delta_{h,n}^*\Delta_{h,n}=|\delta(h,n)|^2I_4\). This identification
is specific to the displayed symbol differences; the general Duhamel
lemma is an operator-norm estimate. Replacing \(|t|\) by \(T\) gives the
compact-time bound. Its right side tends to zero as \(h\to0\).
\(\square\)

For PH10, \(C(n_0)=0\), \(C(n_A)=(2\pi)^3/6\), and
\(C(n_D)=\sqrt3(2\pi)^3/6\), so

\[
\left(\sum_{n\in S}C(n)^2\|a_n\|_2^2\right)^{1/2}
=\frac{(2\pi)^3}{6}=\frac{4\pi^3}{3}.
\]

At its spacing, put \(d=2\pi-6>0\). The axis residual has magnitude
\(d\), the diagonal residual has magnitude \(\sqrt3d\), and their
squared amplitude weights are \(1/4,1/4\). Hence the exact middle
coefficient is

\[
\left(\tfrac14d^2+\tfrac14\,3d^2\right)^{1/2}=d,
\]

and the worked bound is

\[
\|f_{1/12}(t)-f(t)\|_{L^2}
\le(2\pi-6)|t|\le\frac{\pi^3}{108}|t|.
\]

The triangle inequality also gives \(\|f_h(t)-f(t)\|_{L^2}\le2\) for
this normalized packet, so one may use the smaller bound. Neither says
that a fixed nonzero spacing has small error for arbitrarily large time.
Nor does the time bound remove the cutoff: its coefficient contains high
powers of mode momenta. For arbitrary \(L^2\) data, the Fourier-tail argument in the following part
provides convergence without this uniform quadratic coefficient.

\hypertarget{probability-hamiltonian-expectation-and-physical-scope}{%
\section{Probability, Hamiltonian expectation, and physical
scope}\label{probability-hamiltonian-expectation-and-physical-scope}}

Let \(H_S,H_{h,S}\) denote the finite block diagonal Hamiltonians. In
coefficient coordinates, squared norm and continuum Hamiltonian
expectation are

\[
N(a)=\sum_{n\in S}\|a_n\|_2^2,\qquad
\mathcal E(a)=\sum_{n\in S}\langle a_n,H_na_n\rangle,
\]

with centered expectation \(\mathcal E_h(a)\) obtained by replacing
\(H_n\) with \(H_{h,n}\). Hermiticity makes each expectation real: its
complex conjugate is
\(\langle H_na_n,a_n\rangle=\langle a_n,H_na_n\rangle\).

\needspace{4\baselineskip}
\textbf{Proposition PH12 (separate conservation of Hamiltonian
expectation).} Each exact packet evolution preserves the expectation of
its own time-independent Hamiltonian:

\[
\mathcal E(U_S(t)a)=\mathcal E(a),\qquad
\mathcal E_h(U_{h,S}(t)a)=\mathcal E_h(a).
\]

Here the packet operators act on coefficients through the isometry
\(\mathcal K_S\cong\mathcal F_S\). Squared-norm conservation does not
identify either expectation with \(N(a)\). At common initial data,

\[
|\mathcal E_h(a)-\mathcal E(a)|
\le h^2\sum_{n\in S}C(n)\|a_n\|_2^2.
\]

\emph{Proof.} The exponential of a fixed matrix commutes with that
matrix, term by term in its power series. Unitarity gives on each mode

\[
\langle U_H(t)a,H U_H(t)a\rangle
=\langle a,U_H(t)^*HU_H(t)a\rangle
=\langle a,Ha\rangle.
\]

Sum to prove both conservation laws. For common initial coefficients,
Cauchy--Schwarz and PH07 give

\[
|\langle a_n,(H_{h,n}-H_n)a_n\rangle|
\le\|a_n\|_2\|\Delta_{h,n}a_n\|_2
\le h^2C(n)\|a_n\|_2^2.
\]

The finite-sum triangle inequality proves the comparison. \(\square\)

A unit vector in the negative-energy range of \(P_-\) has squared norm
one but expectation \(-E\); a unit vector in the positive-energy range
has expectation \(E\). Norm and energy are therefore different even for
one mode. For \(e_0\), the matrix gives
\(\langle e_0,H(m,p)e_0\rangle=m\), since the kinetic terms have zero
diagonal entry in that component. PH10 has expectation \(1\) for both
Hamiltonians. Its values \(\sqrt{37},\sqrt{109}\) are shell magnitudes,
and its spinors contain both energy signs. Their expectations need not
equal a positive shell value.

The analysis is free and linear. No gauge field, forcing, interaction,
or time-dependent Hamiltonian occurs in these symbols. Other project
modules treat different transport structures; importing them does not
convert these statements into theorems about interacting fields. The
centered propagator is the exact exponential of a modified spatial
symbol. A numerical approximation to that exponential must be checked separately
for its own norm error.

